% ============================================================================
%  Controladores Lógicos Programáveis — apostila do IFF
%  Template-raiz. Copie para apostila/apostila.tex e ajuste edicao/revisao.
%  Compilar com scripts/build.sh (pdflatex + makeindex; latexmk nao existe aqui).
% ============================================================================
\documentclass[11pt,a4paper,oneside]{book}

% --- Codificação e idioma (pdflatex) ---
\usepackage[T1]{fontenc}
\usepackage[utf8]{inputenc}
\usepackage[brazilian]{babel}
\usepackage{lmodern}
\usepackage{microtype}

% --- Cores da identidade (antes do hyperref, que as usa) ---
\usepackage{xcolor}
\definecolor{iffazul}{RGB}{0,84,140}
\definecolor{iffverde}{RGB}{0,122,86}
\definecolor{iffalerta}{RGB}{178,34,34}

% --- Layout ---
\usepackage[a4paper,top=3cm,bottom=2.5cm,left=3cm,right=2.5cm]{geometry}
\usepackage{fancyhdr}
\usepackage{titlesec}
\usepackage{tocloft}
\usepackage{enumitem}
\usepackage{csquotes}

% --- Matemática, tabelas e figuras ---
\usepackage{amsmath,amssymb}
\usepackage{booktabs}
\usepackage{longtable}
\usepackage{array}
\usepackage{graphicx}
\graphicspath{{figuras/}}
\usepackage{adjustbox}
\usepackage{caption}
\usepackage{float}
\usepackage{tikz}
\usepackage{circuitikz}
\usepackage{pdfpages}

% --- Código e caixas didáticas ---
\usepackage{listings}
\usepackage[most]{tcolorbox}
\tcbuselibrary{breakable,skins}

% --- Índices e links (hiperreferências) ---
\usepackage{imakeidx}
\usepackage[
  bookmarks=true,
  bookmarksnumbered=true,
  bookmarksopen=true,
  colorlinks=true,
  linkcolor=iffazul,
  citecolor=iffazul,
  urlcolor=iffazul,
  pdftitle={Controladores Logicos Programaveis},
  pdfauthor={M.Sc. Leonardo de Oliveira Tavares; D.Sc. William da Silva Vianna; M.Sc. Natalia Nogueira Monteiro},
  pdfsubject={Apostila de CLP - Instituto Federal Fluminense}
]{hyperref}

% ============================================================================
%  Identidade e metadados da obra
% ============================================================================
\newcommand{\obraTitulo}{Controladores Lógicos Programáveis}
\newcommand{\obraSubtitulo}{Material didático de automação industrial}
\newcommand{\obraInstituicao}{Instituto Federal Fluminense}
\newcommand{\obraCampus}{Campus Campos dos Goytacazes}
\newcommand{\obraCurso}{Curso Técnico em Automação Industrial}
% Sem \and: esta lista e usada fora de \author/\maketitle.
\newcommand{\obraAutores}{M.Sc. Leonardo de Oliveira Tavares\\%
  D.Sc. William da Silva Vianna\\%
  M.Sc. Natália Nogueira Monteiro}
\newcommand{\obraRevisor}{M.Sc. Karina Terra de Souza Campelo}
\newcommand{\obraEdicao}{3ª edição}
\newcommand{\obraRevisao}{Rev. 06}
\newcommand{\obraData}{Setembro de 2026}
\newcommand{\obraCidade}{Campos dos Goytacazes -- RJ}

\hypersetup{pdftitle={\obraTitulo\ -- \obraRevisao}}

% --- Cabeçalho e rodapé ---
\pagestyle{fancy}
\fancyhf{}
\fancyhead[L]{\small\itshape\nouppercase{\leftmark}}
\fancyhead[R]{\small\thepage}
\renewcommand{\headrulewidth}{0.4pt}
\fancypagestyle{plain}{\fancyhf{}\fancyfoot[C]{\small\thepage}\renewcommand{\headrulewidth}{0pt}}

\titleformat{\chapter}[display]
  {\normalfont\huge\bfseries\color{iffazul}}{\chaptername\ \thechapter}{12pt}{\Huge}
\titlespacing*{\chapter}{0pt}{-20pt}{28pt}

% ============================================================================
%  Estilo de listagem de código de CLP
% ============================================================================
\lstdefinestyle{clp}{
  basicstyle=\ttfamily\footnotesize,
  numbers=left,
  numberstyle=\tiny\color{gray},
  stepnumber=1,
  frame=single,
  rulecolor=\color{gray!50},
  breaklines=true,
  showstringspaces=false,
  tabsize=2,
  captionpos=b,
  extendedchars=true,
  inputencoding=utf8,
  literate=%
    {á}{{\'a}}1 {à}{{\`a}}1 {â}{{\^a}}1 {ã}{{\~a}}1 {ä}{{\"a}}1
    {é}{{\'e}}1 {ê}{{\^e}}1 {è}{{\`e}}1
    {í}{{\'i}}1 {î}{{\^i}}1
    {ó}{{\'o}}1 {ô}{{\^o}}1 {õ}{{\~o}}1 {ò}{{\`o}}1
    {ú}{{\'u}}1 {ü}{{\"u}}1
    {ç}{{\c c}}1 {Ç}{{\c C}}1
    {Á}{{\'A}}1 {À}{{\`A}}1 {Â}{{\^A}}1 {Ã}{{\~A}}1
    {É}{{\'E}}1 {Ê}{{\^E}}1
    {Í}{{\'I}}1 {Ó}{{\'O}}1 {Ô}{{\^O}}1 {Õ}{{\~O}}1 {Ú}{{\'U}}1
    {º}{{\textordmasculine}}1 {ª}{{\textordfeminine}}1
}
\lstset{style=clp}

% ============================================================================
%  Caixas didáticas (usar em todos os capítulos)
% ============================================================================
\newtcolorbox{objetivos}[1][]{
  breakable, colback=iffazul!5, colframe=iffazul, title=Objetivos do capítulo,
  fonttitle=\bfseries, #1}
\newtcolorbox{pratica}[1][]{
  breakable, colback=iffverde!5, colframe=iffverde, title=Na prática,
  fonttitle=\bfseries, #1}
\newtcolorbox{atencao}[1][]{
  breakable, colback=iffalerta!5, colframe=iffalerta, title=Atenção,
  fonttitle=\bfseries, #1}
\newtcolorbox{exercicios}[1][]{
  breakable, colback=gray!5, colframe=gray!60, title=Exercícios,
  fonttitle=\bfseries, #1}

% ============================================================================
%  Capítulos: arquivo ausente gera aviso, não erro
%  (grepar "SKILL-apostila" no log mostra o que ainda falta escrever)
% ============================================================================
\newcounter{capfalta}
\newcommand{\chapterfile}[1]{%
  \IfFileExists{capitulos/#1.tex}%
    {\input{capitulos/#1}}%
    {\stepcounter{capfalta}\typeout{SKILL-apostila: capitulo ausente -> capitulos/#1.tex}}%
}
\AtEndDocument{%
  \ifnum\value{capfalta}>0
    \typeout{SKILL-apostila: \arabic{capfalta} capitulo(s) ainda nao escritos}%
  \fi}

\makeindex[title=Índice remissivo, intoc]

% ============================================================================
\begin{document}

% --- Capa ---
\begin{titlepage}
  \centering
  \vspace*{1cm}
  {\large\bfseries \obraInstituicao\par}
  {\small \obraCampus\par}
  \vspace{1.2cm}
  \rule{\linewidth}{0.5pt}\par
  \vspace{0.8cm}
  {\Huge\bfseries\color{iffazul} \obraTitulo\par}
  \vspace{0.4cm}
  {\large \obraSubtitulo\par}
  \vspace{0.8cm}
  \rule{\linewidth}{0.5pt}\par
  \vfill
  {\large \obraAutores\par}
  \vspace{0.4cm}
  {\small Revisão técnica: \obraRevisor\par}
  \vspace{1.2cm}
  {\obraCurso\par}
  \vspace{0.3cm}
  {\obraEdicao\ -- \obraRevisao\par}
  {\obraData\par}
  \vfill
  {\obraCidade\par}
\end{titlepage}

% --- Folha de rosto ---
\thispagestyle{empty}
\begin{center}
  {\large\bfseries \obraInstituicao}\\[0.2cm]
  {\small \obraCampus}\\[2cm]
  {\Large\bfseries \obraTitulo}\\[0.5cm]
  {\small Material didático elaborado para o \obraCurso.}\\[2cm]
\end{center}
\begin{flushright}
  \begin{minipage}{0.55\textwidth}
    \small Autores:\\ \obraAutores\\[0.6cm]
    Revisão técnica:\\ \obraRevisor\\[0.8cm]
    \obraEdicao\ --- \obraRevisao\\
    \obraData
  \end{minipage}
\end{flushright}
\vfill
\clearpage

% --- Sumário e listas (índice interativo: links clicáveis + bookmarks) ---
\pdfbookmark[0]{Sumário}{sumario}
\tableofcontents
\clearpage
\pdfbookmark[0]{Lista de figuras}{lof}
\listoffigures
\pdfbookmark[0]{Lista de tabelas}{lot}
\listoftables
\clearpage

% ============================================================================
%  PARTE I — FUNDAMENTOS E PROJETO
% ============================================================================
\part{Fundamentos e projeto}
\chapterfile{01-introducao}
\chapterfile{02-evolucao-clps}
\chapterfile{03-estrutura-automacao}
\chapterfile{04-projeto-automacao}

% ============================================================================
%  PARTE II — HARDWARE E FUNCIONAMENTO
% ============================================================================
\part{Hardware e funcionamento}
\chapterfile{05-definicoes-caracteristicas}
\chapterfile{06-componentes-memoria}
\chapterfile{07-cartoes-es}
\chapterfile{08-porte-especificacao}
\chapterfile{09-principio-funcionamento}

% ============================================================================
%  PARTE III — PROGRAMAÇÃO E NORMAS
% ============================================================================
\part{Programação e normas}
\chapterfile{10-linguagens-programacao}
\chapterfile{11-iec61131-3-modelo}
\chapterfile{12-familia-iec61131}
\chapterfile{13-iec-versus-legado}
\chapterfile{14-boas-praticas}
\chapterfile{15-documentacao}
\chapterfile{16-simulacao-teste}

% ============================================================================
%  PARTE IV — INTEGRAÇÃO E SISTEMAS
% ============================================================================
\part{Integração e sistemas}
\chapterfile{17-redes-scada}
\chapterfile{18-iiot-opcua-mqtt}
\chapterfile{19-dados-tempo-real}
\chapterfile{20-interoperabilidade}
\chapterfile{21-convergencia-ti-to}
\chapterfile{22-seguranca-cibernetica}
\chapterfile{23-eficiencia-energetica}

% ============================================================================
%  PARTE V — OPERAÇÃO E CARREIRA
% ============================================================================
\part{Operação e carreira}
\chapterfile{24-manutencao}
\chapterfile{25-formacao-continua}

% ============================================================================
%  APÊNDICES
% ============================================================================
\appendix
\part{Apêndices}
\chapterfile{apA-praticas-laboratorio}
\chapterfile{apB-legado-para-iec}
\chapterfile{apC-glossario}
\chapterfile{apD-referencias}

\printindex

\end{document}
